\section{Constraint Generation for Theorem Proving}\label{sec:theorem_prover}

Such constraints are then used to build class invariant predicates,
stating that there exists an upper bound to the size of the containers
such that, if the containers are bounded before the call to the API,
they are bounded also after the call to the API.

For instance, for \<Ponzi1> the predicate is:
%
\[
  \exists k.\left(\begin{array}{c}
    (\mathit{true}\wedge\<investors>'=1)\Rightarrow\<investors>'\le k\\
    \wedge\\
    (\<investors>\le k\wedge\<investors'>=\<investors>+1)\Rightarrow\<investors>'\le k
  \end{array}\right)
\]
%
that a theorem prover should reject or fail to prove that it holds. Therefore,
the analyzer would issue a boundedness warning.

For \<Ponzi2>, the predicate is:
%
\[
  \exists k.\left(\begin{array}{c}
    (\mathit{true}\wedge\<investors>'=1)\Rightarrow\<investors>'\le k\\
    \wedge\\
    (\<investors>\le k\wedge\<investors>'=\<investors>+1)\Rightarrow\<investors>'\le k\\
    \wedge\\
    (\<investors>\le k\wedge\<investors>'=\<investors>)\Rightarrow\<investors>'\le k\\
  \end{array}\right)
\]
%
that a theorem prover should reject or fail to prove that it holds. Therefore,
the analyzer would issue a boundedness warning.

For \<Ponzi3>, the predicate is:
%
\[
  \exists k.\left(\begin{array}{c}
    (\mathit{true}\wedge\<investors>'=1)\Rightarrow\<investors>'\le k\\
    \wedge\\
    (\<investors>\le k\wedge\<investors><100\wedge\<investors>'=\<investors>+1)\\
    \Rightarrow\<investors>'\le k
  \end{array}\right)
\]
%
that a theorem prover should be able to prove true, for instance by taking
$k=100$ (or $k=200$\ldots). Therefore, the boundedness analyzer would not issue
any warning.

For \<Lottery1>, there is a predicate for \<rounds> and another for
\<rounds[any].sales>. The former is irrelevant, since \<rounds> is not
iterated. The latter is (by writing \<ras> for \<rounds[any].sales>):
%
\[
  \exists k.\left(\begin{array}{c}
    (\mathit{true}\wedge\<ras>'=0)\Rightarrow\<ras>'\le k\\
    \wedge\\
    (\<ras>\le k\wedge\<ras>\le\<ras>'\le\<ras>+1)\Rightarrow\<ras>'\le k\\
    \wedge\\
    (\<ras>\le k\wedge\<ras>'=\<ras>)\Rightarrow\<ras>'\le k
  \end{array}\right)
\]
%
that a theorem prover should reject or fail to prove that it holds. Therefore,
the analyzer would issue a boundedness warning.

For \<Lottery2>, the predicate for \<rounds[any].sales> is:
%
\[
  \exists k.\left(\begin{array}{c}
    (\mathit{true}\wedge\<ras>'=0)\Rightarrow\<ras>'\le k\\
    \wedge\\
    (\<ras>\le k\wedge\<ras>\le\<ras>'\le\max(5000,\<ras>))\Rightarrow\<ras>'\le k\\
    \wedge\\
    (\<ras>\le k\wedge\<ras>'=\<ras>)\Rightarrow\<ras>'\le k
  \end{array}\right)
\]
%
that a theorem prover should be able to prove true, for instance by taking
$k=5000$ (or $k=8000$\ldots). Therefore, the boundedness analyzer would not
issue any warning.

\textbf{Which theorem prover is able to deal with such constraints?}
